\chapter{Using FsYacc as a parser generator \label{chapter:ply}}
FsYacc generates a bottom-up LALR parser from an \texttt{.fsy} grammar. Its semantic actions are F\# expressions.
FsLex supplies tokens with the same generated token type. The scanner, generated parser and domain evaluator are
separate modules; a small frontend composes them with one lexical buffer.

\section{A symbolic calculator}
The original calculator grammar remains the mathematical specification.
@@figure:ply-yacc-interpreter:fig:calculator.g@@

The syntax tree records numbers, variable names, arithmetic operators and statements with explicit union cases.
The scanner and parser use that domain representation as follows.
@@snippet:calculator-types@@
@@snippet:calculator-lexer@@
@@snippet:calculator-parser@@

\texttt{\%token <float> NUMBER} declares a token with a floating-point payload.
\texttt{\%start} selects the entry nonterminal, and \texttt{\%type} fixes its semantic result type.
In an action, \texttt{\$1} is the first right-hand-side value, \texttt{\$2} the second, and so on.
The number action constructs \texttt{Number \$1}; no unchecked numeric tuple index is carried into evaluation.
The grammar's separate product and factor nonterminals encode arithmetic precedence structurally.

The environment is an immutable map from names to values. Assignment returns a new environment;
a lookup either finds a value or reports an undefined variable. The evaluator therefore has an explicit state boundary.
@@snippet:calculator-evaluate@@
\href{../../Chapter-07/01-Calculator.ipynb}{\nolinkurl{Chapter-07/01-Calculator.ipynb}}
contains the full generated-frontend example and a separately callable interactive loop.
Running all demonstration cells never blocks waiting for terminal input.

\section{Symbolic differentiation}
To differentiate we retain the expression tree rather than evaluating it to a number.
Pattern matching implements the sum, product, quotient, chain and power rules directly over the recursive type.
@@snippet:differentiate@@
The differentiation notebook checks its symbolic result against a numerical finite-difference approximation.
Its pretty-printer retains the parentheses needed to parse the tree again.
See \href{../../Chapter-07/02-Differentiator.ipynb}{\nolinkurl{Chapter-07/02-Differentiator.ipynb}}.

\section{Implementing a simple interpreter \label{chapter:interpreter}}
The source language has assignments, blocks, conditionals, loops and calls to predefined functions.
This source language is distinct from the F\# program that interprets it.
@@figure:ply-yacc-interpreter:fig:Pure.g@@
@@figure:ply-yacc-interpreter:fig:sum.sl@@
The example computes $\sum_{i=1}^{n^2} i$ and prints the result.
Its syntax tree uses \texttt{Block}, \texttt{Assign}, \texttt{If} and \texttt{While} cases;
there is no nested tuple whose first string must be interpreted as an instruction name.
@@snippet:interpreter-types@@
@@snippet:interpreter-parser@@

Evaluation receives a map of typed values and an explicit I/O record.
Integer literals, input, addition, subtraction and multiplication retain arbitrary-precision integers.
The simple language uses true division, which returns a floating-point value; modulo follows Python's
floor-quotient rule, including negative operands. Mixed integer/real comparisons preserve integer precision.
That record is the actual boundary between deterministic evaluation and reading or printing.
Undefined variables and unsupported constructs produce diagnostics rather than fabricated values.
@@snippet:interpreter-evaluate@@

Execution returns the updated map. A block folds over its statements; a loop recursively threads its current
map through its body. This makes the loop invariant visible in both the type and the control flow.
@@snippet:interpreter-execute@@
@@figure:ply-yacc-interpreter:fig:factorial.sl@@
\href{../../Chapter-07/03-Interpreter.ipynb}{\nolinkurl{Chapter-07/03-Interpreter.ipynb}}
contains the scanner/parser construction and the simple sum program.
The factorial source example requires the extended language with function definitions.
The corresponding \href{../../Chapter-07/Interpreter-Frame.ipynb}{\nolinkurl{Chapter-07/Interpreter-Frame.ipynb}}
adds function frames and return behavior through explicit types.

\section{Check your understanding}
\begin{enumerate}
\item Explain the roles of scanner, generated parser, syntax tree and evaluator.
\item Why does a typed semantic action reject a malformed AST before the notebook executes it?
\item Extend the simple interpreter with a comparison or a statement and update every affected match.
\item Explain which effects require I/O and which transformations can operate solely on syntax trees.
\end{enumerate}
